\section{Problem: Identifying Faithfulness}

\paragraph{Local relevance.}
Let $q$ be a query and $P=\{p_1,\ldots,p_k\}$ a retrieved context set. A
standard local query-passage score measures how much each passage resembles
the query, often through embedding cosine similarity. We write the mean local
similarity as
\[
  \simq(P)=\frac{1}{k}\sum_i \cos(e(q),e(p_i)).
\]
This score is useful for retrieval, but it treats the passages as separable
items.

\paragraph{Context coherence.}
The Context Coherence Score is a set-level statistic:
\[
  \ccs(P)=\frac{2}{k(k-1)}\sum_{i<j}\cos(e(p_i),e(p_j)).
\]
It measures whether retrieved passages are mutually similar. CCS can expose
structure that $\simq$ discards, but it is not assumed to be a sufficient
condition for answer support.

\paragraph{Faithfulness.}
A generator $G$ produces $a=G(q,P)$. We use faithfulness to mean answer support
by the retrieved context: the answer should be entailed by, or at least
directly grounded in, $P$. This is a property of the set and the answer
together. It is not implied by high local similarity, because a context set can
look query-relevant while omitting the decisive evidence, duplicating the same
partial fact, or combining passages that invite an unsupported synthesis.
High CCS has the same problem in another form: topically coherent passages may
all be answer-absent.

\paragraph{Identifiability.}
We use \emph{identifiability} in the empirical sense. A reported statistic
identifies a faithfulness claim only if the claimed direction would remain
distinguishable after the relevant alternatives are separated. If a method has
higher $\simq(P)$ but lower support, local similarity did not identify
faithfulness. If a coherence effect disappears after matching $\simq(P)$,
coherence alone did not identify the effect. If different automatic judges
give different intervention magnitudes on identical generations, a single
metric did not identify the effect size.

\paragraph{Research hypotheses.}
The controlled suite tests five hypotheses that are often implicit in RAG
papers:
\begin{description}[leftmargin=1.2em,itemsep=2pt,topsep=2pt]
  \item[H1: Coherence causality.] At fixed $\simq(P)$, higher CCS increases
  answer faithfulness.
  \item[H2: Scale robustness.] The refinement paradox remains large and
  statistically stable at larger $n$ and multiple seeds.
  \item[H3: Metric invariance.] Faithfulness conclusions are stable across
  reasonable automatic evaluators.
  \item[H4: Threshold transfer.] A refinement policy tuned on one dataset
  transfers to others without large recovery loss.
  \item[H5: Method dominance.] A coherence-aware policy dominates alternative
  RAG interventions after latency and indexing cost are included.
\end{description}
The results reject H1, weaken H2, reject H3, qualify H4 as dataset-dependent,
and reject H5. The remaining positive claim is narrower: CCS is a diagnostic
variable that can reveal structure missed by local similarity, especially in
the coherent-noise control, but it is not a validated mechanism.
